\chapter*{ВВЕДЕНИЕ}
\addcontentsline{toc}{chapter}{ВВЕДЕНИЕ}

В соответствии с~\cite{international_federation_of_robotics}, сервисные роботы для потребительского использования продемонстрировали устойчивый рост на 11 \%, и в 2024 году было продано около 20,1 млн единиц.

Согласно~\cite{ronzhin2006speech}, повышению эффективности ЧМВ способствует использование наиболее естественных для человека средств коммуникации. 
Здесь наблюдается, с одной стороны, перенос людьми своих ожиданий с межличностной коммуникации на ЧМВ и, с другой стороны, использование мультимодальной коммуникации, в том числе аффективной.
При этом исследования в области ЧМВ показывают, что люди при взаимодействии с роботами используют социальные и коммуникативные механизмы, характерные для межличностного взаимодействия~\cite{park2012lawofattraction}. 

Опираясь на~\cite{review_of_personality_of_hri}, можно сказать, что люди положительно воспринимают эффект персонализации роботов-собеседников, в частности, предпочитают более экстравертных роботов.

Ограничение данной работы --- это применимость ДС в рамках развлекательной, образовательной и социальной функциональности согласно с видами функциональности ДС, описанными в~\cite{Volkova2023_dialogue_systems}.

Таким образом, решение задачи персонализации роботов-собеседников является актуальной, поскольку позволяет повысить естественность ЧМВ и качество восприятия робота.
Перспективной представляется модификация метода ведения диалога роботом Ф-2 на основе учёта персональных черт, назначаемых роботу-собеседнику.

Цель работы -- разработка модифицированного метода ведения диалога роботом Ф-2 с учётом персональных черт робота.

Для достижения поставленной цели необходимо решить следующие задачи:

\begin{itemize}
	\item провести анализ существующих подходов к описанию персональных черт;

	\item провести обзор существующих решений (ДС), в которых применены различные ПОПЧ;

	\item провести анализ исходного метода ведения диалога роботом Ф-2;
	
	\item разработать модифицированный метод ведения диалога роботом Ф-2 с учётом персональных черт робота;
	
	\item реализовать разработанный метод;
	
	\item провести оценку качества результатов реализованного метода.

\end{itemize}
